% ---------------------------------------------------------------
\section{Results}
\label{sec:results}
% ---------------------------------------------------------------

\begin{table}[htbp]
\centering
\caption{Parameters and values explored. Baseline values in italics.}
\label{tab:paramvalues}
\begin{tabular}{lllp{7.5cm}}
\toprule
Symbol & Name & Baseline & Values explored \\
\midrule
\multicolumn{4}{@{}l@{}}{\textit{Corpus construction (see Tables~\ref{tab:params} and~\ref{tab:windows})}} \\[2pt]
$t_x$          & Time window            & \textit{5}              & 1,\ 2,\ 3,\ 4,\ \textit{5},\ 6 \\
$F$            & Field subset           & \textit{$A$}            & $E$,\ $B$,\ \textit{$A$} \\
$\tau_U$       & Institution threshold  & \textit{10}             & \textit{10},\ 8 \\[4pt]
\multicolumn{4}{@{}l@{}}{\textit{Matrix construction}} \\[2pt]
$\rho$         & Reference normalisation & \textit{0 (fixed)}     & \textit{0} (fixed),\ 1 (full) \\
$\mathbf{m}$   & Block mask             & \textit{$(0,1,1,0)$}    & $(1,0,0,0)$,\ $(0,0,0,1)$,\ \textit{$(0,1,1,0)$},\ $(1,1,1,1)$ \\
$\chi$         & Mixing weight          & \textit{0.50}           & \textit{0.50},\ $\chi^*=N_u/(N_s{+}N_u)\approx[FILL]$ \\[4pt]
\multicolumn{4}{@{}l@{}}{\textit{Ranking}} \\[2pt]
$\alpha$       & Damping factor         & \textit{0.85}           & 0.50,\ \textit{0.85} \\
\bottomrule
\end{tabular}
\end{table}

Table~\ref{tab:paramvalues} summarises the parameter space explored.
The baseline is the bipartite SI/IS mode ($\mathbf{m}=(0,1,1,0)$,
$F=A$, $t_x=5$, $\tau_U=10$, $\rho=0$, $\alpha=0.85$).
Results are organised around four analytical questions:
(\textit{i}) what is the source-mode community structure and how do
Economics and Business journals rank within it;
(\textit{ii}) how does the institution network dissolve that community
structure;
(\textit{iii}) what is the primary ranking under the bipartite baseline;
and (\textit{iv}) how do the full-joint mode and parameter sensitivity
affect these findings.
Parts~\ref{sec:results_ss}--\ref{sec:results_sensitivity} address these
in turn.

% ---------------------------------------------------------------
\subsection{Source ranking and community structure: the SS mode}
\label{sec:results_ss}
% ---------------------------------------------------------------

We begin with the source-only ($\mathbf{m}=(1,0,0,0)$) case, which
restricts attention flow to the journal--journal channel and produces a
ranking on $\mathfrak{S}$ alone. This case is analytically the simplest and
provides a reference for the community structure that the subsequent
bipartite and full-joint modes partially dissolve.

\paragraph{SS baseline ranking ($F=A$).}
The prestige-per-work scores $v_s$ span [FILL: order-of-magnitude range, e.g.\ two orders of magnitude], with a
few journals at $v_s \gg 1$ and a long tail below the corpus
average. Table~S1 lists the top [FILL: $N$] sources by $v_s$ under
$F=A$; Figure~\ref{fig:v_rank} (top panel) shows the full distribution.
[FILL: 1--2 sentences describing which journals appear at the top and
whether this matches prior expectations from impact-factor rankings.]

\paragraph{Economics vs.\ Business subsets.}
Restricting to $F=E$ (Economics sources only) and $F=B$ (Business
sources only) yields separate rankings on the restricted corpora.
Tables~S2 and~S3 list the top sources under each restriction.
The most striking movement between the $F=A$ and $F=E$ rankings involves
journals that span the Economics--Finance--Business boundary: under $F=E$,
intra-corpus references from Business sources are removed, cutting the
prestige flow that these bridge journals receive from the Business
community. [FILL: specific example journals and rank shifts.]
The corresponding rank shifts are shown in Figure~[FILL].

\paragraph{Community structure: second eigenpair of $\mathbf{H}_{SS}^\top$.}
The Katz--Hubbell resolvent has the spectral expansion
$\boldsymbol{\pi} = (1-\alpha)\sum_i [v_i^\top\boldsymbol{\mu}/(1-\alpha\lambda_i)]\,u_i$,
where $(u_i,v_i,\lambda_i)$ are the right eigenvectors, left eigenvectors,
and eigenvalues of $\mathbf{H}_{SS}^\top$.
The $i=1$ term recovers $\boldsymbol{\pi}$ itself ($\lambda_1=1$);
the $i=2$ term is the dominant correction, amplified by $1/(1-\alpha\lambda_2)$.
The second left eigenvector $\boldsymbol{\varphi}_2$ of $\mathbf{H}_{SS}^\top$
therefore encodes the dominant partition of sources beyond the uniform
Katz ranking.

Under $F=A$, the second eigenvalue is $\lambda_2^{SS}=[FILL]$,
giving spectral gap $g^{SS}=1-\lambda_2^{SS}=[FILL]$ and community
amplification $A^{SS}=1/(1-\alpha\lambda_2^{SS})=[FILL]$.
Figure~\ref{fig:community} (Panel~A) plots the elements of
$\boldsymbol{\varphi}_2$ sorted by value, coloured by field label
($F=E$: blue; $F=B$: orange).
[FILL: describe the sign partition -- expected to be a clean E/B split
with finance journals near zero.]
The sign split of $\boldsymbol{\varphi}_2$ directly identifies the
Economics and Business communities from the citation structure alone,
with no external classification imposed.
Finance journals (e.g.\ the \textit{Journal of Finance},
\textit{Journal of Financial Economics}, \textit{Review of Financial Studies})
appear near $\varphi_{2,s}\approx 0$, confirming their role as citation
bridges between the two communities: their prestige under $F=A$ derives
from both sides of the sign partition, explaining the rank drops observed
in the $F=E$ restricted ranking.

Table~S4 presents the community partition table for all sources, sorted
by $\varphi_{2,s}$, with $F$-label, $v_s$, and rank by $v_s$ under
the SS baseline.

% ---------------------------------------------------------------
\subsection{Institution ranking and community diffusion: the II mode}
\label{sec:results_ii}
% ---------------------------------------------------------------

The institution-only mode ($\mathbf{m}=(0,0,0,1)$) restricts attention
flow to the institution--institution channel.
Since authors at a single institution typically span both Economics and
Business, the II citation network is structurally more integrated than
the journal network.

\paragraph{II baseline ranking ($F=A$).}
Table~S5 lists the top [FILL: $N$] institutions by $v_u$ under II.
[FILL: describe leading institutions -- expected to be large, multi-faculty
research universities spanning both fields. Note any contrast with
single-field specialist institutions.]

\paragraph{Community structure: second eigenpair of $\mathbf{H}_{II}^\top$.}
Under $F=A$, the second eigenvalue is $\lambda_2^{II}=[FILL]$,
giving spectral gap $g^{II}=[FILL]$ and amplification
$A^{II}=[FILL]$.
As expected, $g^{II}>g^{SS}$: the institution network mixes the
Economics and Business communities more rapidly than the journal network,
because institutions co-locate researchers from both fields.
Figure~\ref{fig:community} (Panel~B) shows $\boldsymbol{\varphi}_2$
for institutions; the distribution is [FILL: expected to be noisier, with
more units near zero].
The smaller amplification $A^{II}<A^{SS}$ implies that community
membership has less influence on institution rankings than on source
rankings: prestige is more evenly distributed across the institution
network.

% ---------------------------------------------------------------
\subsection{Bipartite baseline: sources and institutions ranked jointly}
\label{sec:results_bipartite}
% ---------------------------------------------------------------

The bipartite SI/IS mode ($\mathbf{m}=(0,1,1,0)$) is the paper's
primary specification.
Citation attention flows from sources to institutions (via
$\mathbf{H}_{SI}$) and from institutions back to sources (via
$\mathbf{H}_{IS}$), so each reference traverses one S$\to$I step and one
I$\to$S step.
Under the $\sqrt{\alpha}$ per-step convention (Section~\ref{sec:networks}),
the round-trip attenuation per reference equals $\alpha$, matching the
single-step attenuation in the SS and II modes and making all four modes
directly comparable.

\paragraph{Source ranking.}
Table~1 lists the top [FILL: $N$] sources by $v_s$ under the bipartite
baseline.
[FILL: describe leading journals and compare qualitatively to SS.
Expected: finance bridge journals rise relative to SS because their
institutional affiliates span both communities, giving them indirect
access to prestige from both the Economics and Business sides.]
The rank correlation between the bipartite and SS source rankings is
$\rho_s=[FILL]$ (Spearman), indicating [FILL: substantial stability /
meaningful reordering].

\paragraph{Institution ranking.}
Table~2 lists the top [FILL: $N$] institutions by $v_u$ under the
bipartite baseline.
[FILL: describe leading institutions. Note that bipartite institution
scores reflect the institutions' role as citation intermediaries between
sources, not direct institution-to-institution citation flow.]
The rank correlation between bipartite and II institution rankings is
$\rho_u=[FILL]$.

\paragraph{Community structure: second eigenpair of $\mathbf{M}_S^\top$.}
In the bipartite mode, the relevant community operator is the one-mode
projection $\mathbf{M}_S=\mathbf{H}_{SI}\mathbf{H}_{IS}$.
Its second eigenvalue is $\lambda_2^{M_S}=[FILL]$, giving spectral gap
$g^{\mathrm{bip}}=[FILL]$ and amplification
$A^{\mathrm{bip}}=1/(1-\alpha\lambda_2^{M_S})=[FILL]$.
The ordering $g^{SS}<g^{\mathrm{bip}}<g^{II}$ [FILL: confirm or note
if violated] reflects the intermediate mixing role of the bipartite walk:
routing citation flow through institutions partially blurs the E/B
community boundary, but less completely than the direct II channel
because each source's institutional affiliates are not uniformly
distributed across both communities.

The institution community vector is recovered as
$\boldsymbol{\varphi}_2^I=\mathbf{H}_{SI}^\top\boldsymbol{\varphi}_2^S$
(normalised), providing an institution-side partition induced by the
source walk. [FILL: brief description of whether this aligns with II
community partition or differs.]

% ---------------------------------------------------------------
\subsection{Full-joint mode and the mixing weight $\chi$}
\label{sec:results_full}
% ---------------------------------------------------------------

Adding the direct diagonal blocks ($\mathbf{m}=(1,1,1,1)$) combines
source-mode and institution-mode citation flow via the mixing weight
$\chi$. We examine two calibrations: the reference case $\chi=0.5$
and the dimensionally neutral calibration
$\chi^*=N_u/(N_s+N_u)\approx[FILL]$, at which the expected
prestige per unit is equal across sources and institutions.

\paragraph{Full joint at $\chi=0.5$.}
Tables~S6 and~S7 list the top sources and institutions under the full-joint
mode at $\chi=0.5$.
The rank correlations against the bipartite baseline are
$\rho_s=[FILL]$ (sources) and $\rho_u=[FILL]$ (institutions).
[FILL: characterise whether adding the diagonal blocks materially
reorders the rankings or leaves them stable.]
The second eigenvalue of the full-joint $\mathbf{H}^\top$ is
$\lambda_2^{\mathrm{full},0.5}=[FILL]$,
with amplification $A^{\mathrm{full},0.5}=[FILL]$.
The second eigenvector $\boldsymbol{\varphi}_2$ (length $N=N_s+N_u$)
encodes [FILL: E/B community split, S\/I mode split, or a combination
-- to be determined from data].

\paragraph{Dimensionally neutral calibration $\chi^*$.}
At $\chi=0.5$ with $N_u>N_s$, institution units receive less prestige
flow per unit than source units. The value $\chi^*\approx[FILL]$
restores balance.
Figure~\ref{fig:community} (Panel~D) scatters the second eigenvector
elements at $\chi=0.5$ against those at $\chi^*$ for sources,
coloured by $F$-label.
[FILL: expected -- minor shifts in magnitude, stable sign partition,
confirming that the E/B community structure is robust to dimensional
rebalancing.]
The rank correlations are $\rho_s=[FILL]$ and $\rho_u=[FILL]$.

\paragraph{Spectral gap summary.}
Figure~\ref{fig:community} (Panel~C) summarises the spectral gap
$g=1-\lambda_2$ and community amplification $A=1/(1-\alpha\lambda_2)$
across all five cases (SS, bipartite $M_S$, II, full $\chi=0.5$,
full $\chi^*$).
The [FILL: expected monotone] progression in $g$ from the journal-only
mode to the institution-only mode quantifies the degree to which
community membership influences rankings in each case.

% ---------------------------------------------------------------
\subsection{Parameter sensitivity}
\label{sec:results_sensitivity}
% ---------------------------------------------------------------

Table~S8 reports Spearman rank correlations between each Stage~1
parameter variant and the bipartite baseline, for sources and
institutions separately.
The main findings are as follows.

\paragraph{Reference normalisation ($\rho$).}
Switching from fixed-count ($\rho=0$, baseline) to full-count ($\rho=1$)
down-weights works with many references, disproportionately affecting
review articles and handbooks. The rank correlation is
$\rho_s^{\rho}=[FILL]$ for sources and $\rho_u^{\rho}=[FILL]$ for
institutions, indicating [FILL: degree of robustness].

\paragraph{Damping factor ($\alpha$).}
At $\alpha=0.5$ (lower than baseline $0.85$) ranking is more heavily
influenced by the prior and less by citation network structure.
The rank correlation with the baseline is $\rho_s^{\alpha}=[FILL]$,
[FILL: comment on whether rankings are stable].

\paragraph{Institution threshold ($\tau_U$).}
Relaxing the institution threshold to $\tau_U=8$ adds [FILL: $\Delta N_u$]
institutions, primarily smaller specialised research units.
The rank correlation for sources is $\rho_s^{\tau}=[FILL]$;
for institutions, the comparison is over the smaller intersection set
retained under both thresholds.

\paragraph{Field restriction ($F$).}
The $F=E$ and $F=B$ sensitivity comparisons for source rankings are
discussed in Part~\ref{sec:results_ss} above as substantive findings
rather than robustness checks; they are included here for completeness.

% ---------------------------------------------------------------
\subsection{Time-series stability}
\label{sec:results_timeseries}
% ---------------------------------------------------------------

The Stage~2 time-series sweep holds all parameters at baseline values
and varies $t_x\in\{1,2,3,4,6\}$.
Figure~\ref{fig:timeseries} shows the rank trajectories of the top
[FILL: $N$] sources across windows.
[FILL: describe whether leading journals are stable across windows or
show systematic trends. Expected: high rank correlation between adjacent
windows; some specialist journals show greater volatility.]
Table~S9 reports pairwise Spearman rank correlations between adjacent
time windows, with correlations of $\rho=[FILL]$ [FILL: or a range]
across the [FILL: $t_x=1$ to $t_x=6$] sequence, confirming [FILL:
stability / instability of the rankings over time].
